\chapter{需求分析与设计原则}
\chapterlead{这些原则大多不是事先想好的，而是在某次出错之后补上的。}

\section{使用场景}

我们设想的使用者是三类人。第一类是\textbf{卖方或买方研究员}，拿到一家新公司的年报后，希望在一天之内得到一份数字可靠的初稿和一个能改假设的估值模型，
再把自己的判断加进去。第二类是\textbf{研究所的质控与合规人员}，需要在发布前核对研报中的数字，确认没有抄错、算错、单位错和期间错。
第三类是\textbf{学习投研的学生和个人投资者}，希望看到一份研报的数字从年报哪一页来、增速怎么算、估值假设是什么。

这三类需求有一个共同点：结论可以讨论，数字不能出错。研究员可以不同意模型对行业的判断，但不会接受一份连营业收入都抄错的初稿；
质控人员需要的是能指出“错在哪一页、属于哪类错误”的核查结果，而不是一句“该数字可能有误”；学生需要的是能一路追溯的计算过程。
因此系统设计的第一优先级是数字可信，第二优先级才是分析深度和行文质量。

\section{竞赛要求与我们的对应}

竞赛对作品有几项硬性要求。表~\ref{tab:pr-req} 列出每项要求在系统中的落实方式。

\begin{table}[htbp]
\centering
\caption{竞赛硬约束与落实方式}
\label{tab:pr-req}
\small
\begin{tabularx}{\linewidth}{@{}>{\raggedright\arraybackslash}p{3.6cm} L@{}}
\toprule
\hd{竞赛要求} & \hd{系统中的落实} \\
\midrule
至少一个大语言模型作为核心推理引擎 & 五个生产角色与盲评质检员均为 Claude 模型会话（生成用 Opus，复核与质检用 Sonnet），负责阅读、抽取、调研、写作、改写与评审 \\
\rowband
输出为结构化数据、表格或报告 & 真值表与分析结果（JSON、Markdown）、估值工作簿（15 个工作表、公式联动）、纠错核查结果（JSON 与报告）、研报 PDF、质检评分（JSON） \\
全程记录文件访问、工具调用与计算过程 & 引擎每次运行写 \file{events.jsonl} 与 \file{manifest.json}，记录每个读入文件和产物的哈希；每个模型会话的完整事件流保存在工作区 \file{.pipeline/agents/} \\
\rowband
提交完整源码、依赖与运行说明 & 全部源码、\file{requirements.txt}、\file{pyproject.toml}、运行说明与样例数据；样例评测一条命令复现 \\
第三方代码如实列明 & \file{THIRD\_PARTY.md} 列出名称、版本、来源、许可证与使用范围；参照的开源估值插件代码不进仓库，自研引擎与其对照关系如实说明 \\
\rowband
估值与投资判断写明假设、适用范围与风险，区分事实、推论与观点 & 估值假设单独成文件并要求写依据；数字护栏把正文每个数字归入事实、推论或观点；研报必须包含假设、口径简化说明与“不构成投资建议” \\
\bottomrule
\end{tabularx}
\end{table}

\paragraph{为什么全部使用 Claude。}
大赛鼓励优先使用国产模型。我们评估后决定生成与质检全部使用 Claude，理由有三点。其一，我们的多智能体架构以无头 Claude Code 作为运行框架（harness），
它原生提供文件读写、命令执行、网络检索、工具白名单和完整的事件流记录，正好满足“隔离工作区、只靠文件交接、全程可追溯”的要求，
自建同等能力的框架会占用大量开发时间。其二，一份年报动辄两三百页，分析师会话要在其中反复翻页、核对、改正，需要长上下文与稳定的工具调用。
其三，系统的可信度并不依赖模型本身：所有计算、校验和拦截都在不调用模型的确定性引擎里完成，模型只是引擎的“使用者”。
引擎与模型之间的接口是文件（\file{facts.json}、\file{narrative.md} 等）和命令行，更换模型只需替换会话启动命令，引擎、闸门和评分标准都不受影响。

\section{设计原则}

下面七条原则写在仓库的开发指南里，也是代码审查时的检查项。每条后面附上它的来由。

\subsection{模型不做算术}
所有数值计算都由确定性 Python 完成，包括增速、占比、利润率、勾稽、预测和估值，连“约三倍”“接近一半”这类约数也不允许模型心算，
必须由程序算出后引用。第 1 章提到的“同比增长 30.45\%”就是心算的产物。这条原则同时简化了核查：
正文里所有派生数字都应该能在本次运行的计算产物中找到，找不到就说明不是程序算的。

\subsection{每个数字都可以追溯}
抽取结果必须带来源 $\sigma = (\text{文件}, \text{页码}, \text{表块编号}, \text{行}, \text{列})$。模型抄的数字必须在所写页码的原文中按数字边界找到，
例如抄“1,688.38”时，原文中的“11,688.38”不算命中。外部资料的摘句必须包含所登记的数字。研报里的外部数字必须带 \code{[S编号]} 指向来源登记表。
这样任何一个数都可以从研报正文一路查回年报的某一页某一格，或者某个网页的某一句。

\subsection{两条通道不一致时不替任何一方做选择}
模型抽取与规则通道的结果不一致时，系统记为问题并阻断估值，交回分析师会话回原文核对。我们考虑过“取规则通道”或“取置信度高的一方”，
但两条通道出错的方式完全不同：规则通道错在表格结构（错列、续页、折行），模型错在阅读（看错行、看错期间），事先无法断定哪一方更可靠。静默选择会把错误藏进真值表，
之后所有估值和研报都建立在它上面。阻断的代价是多一轮核对，这个代价是值得的。

\subsection{只有一张真值表}
财报分析产出的真值表是估值、纠错和研报写作的唯一数据来源。下游模块不允许重新抽取或重新计算基础数字。
年报里“营业收入”和“营业总收入”、“净利润”和“归母净利润”常常同时出现，如果估值和纠错各自取数，两个模块很容易落在不同口径上，
而且彼此都察觉不到。只保留一张表，口径只需要在抽取时判断一次，并且把判断理由写进抽取文件。

\subsection{误报控制优先于召回}
纠错核查时，一个把正确数字判成错误的系统比漏掉一个错误的系统更难用，因为使用者会很快失去信任，转而自己逐条核对。
所以容差按研报数字的展示精度计算（写“58.62 亿元”的人，只对百万元级别负责）；预测性表述和缺少真值的断言判“无法核实”，不计为错误；
错误类型必须找到精确匹配才下结论，否则只说“数值错误”。这一原则也用在了规则通道的表格解析上：宁可漏取，不可错配。

\subsection{会话之间隔离，评分标准对生成方不可见}
研报流水线的每个角色是一次独立会话，不继承上一个角色的上下文，只读工作区里的文件。工作区建在代码仓库之外。
这样做有两个来由。一是模型会“照抄”：仓库里存有别家公司的分析框架和文稿，会话一旦能读到，就倾向于套用现成的结构，
而不是从这家公司的年报出发。二是“应试”：如果生成方能看到评分标准，它会去满足条目字面，而不是把报告写好。因此评分标准和完整的评测结果都不进入工作区，回传的问题单只含理由与证据。

\subsection{不预设行业框架，文风交给编辑}
技能说明（skill）只规定方法与约束，不提供“白酒怎么分析”“光模块看哪些指标”之类的行业模板，驱动因素由模型从年报与资料包中推导。
预先写好的框架等于提前替模型做了分析，而这份分析没有经过这家公司年报的检验；系统要面对任意一家 A 股公司，模板也写不过来。
茅台的推导过程保留在技能说明里，但只作为“如何推导”的示例，明确标注不是模板。对于文风，我们曾考虑用正则表达式拦截“值得注意的是”“不是……而是……”这类句式，
最后放弃了：规则拦截只会让模型换一种同样生硬的说法。现在的做法是设一个独立的复核编辑会话，以没有参与写作的质控编辑的身份通读改写。

\begin{keybox}
{\heiti\bfseries\color{ink}一句话概括}\quad 能用程序判断的，一律交给程序；必须由模型判断的，用隔离、闸门和盲评来约束它。
\end{keybox}
